\documentclass[11pt]{article}
\ifdefined\pdftrailerid\pdftrailerid{}\fi
\usepackage[T1]{fontenc}
\usepackage{lmodern,amsmath,amssymb,amsthm,booktabs}
\usepackage[margin=1in]{geometry}
\usepackage[hidelinks]{hyperref}
\usepackage{microtype}
\hypersetup{pdftitle={Oscillator arrays in pyflock: Mathematical reference}}
\newcommand{\R}{\mathbb R}
\newcommand{\C}{\mathbb C}
\newcommand{\dd}{\mathrm d}
\newcommand{\diag}{\mathop{\mathrm{diag}}}
\newtheorem{proposition}{Proposition}
\title{Oscillator arrays in pyflock\\[3pt]\large Mathematical reference}
\author{}
\date{}
\begin{document}
\maketitle
\vspace{-2em}

Pyflock computes trajectories, spectra and transient observables for linear
oscillator arrays. The principal examples follow Cantos, Hammond and
Veerman~\cite{signal,transients}, the next-nearest-neighbour study of
Herbrych et al.~\cite{nnn}, and Lyons's heterogeneous arrays~\cite{lyons}.
This reference gives the equations and the hypotheses behind the comparisons.
The repository's coverage map identifies the source figure and any changed
parameters or measurement rule for each experiment.

\section{Homogeneous arrays}

Let $z_k(t)$ be displacement from an equally spaced formation. For offsets
$J=\{-2,-1,0,1,2\}$, the interior equation is
\begin{equation}\label{eq:model}
 \ddot z_k=g_x\sum_{j\in J}\rho_{x,j}z_{k+j}
          +g_v\sum_{j\in J}\rho_{v,j}\dot z_{k+j}.
\end{equation}
Decentralization means $\sum_j\rho_{x,j}=\sum_j\rho_{v,j}=0$.
Consequently uniform displacement and velocity, $z_k(t)=a+bt$, solve the
interior equations. The usual normalized convention has
$\rho_{x,0}=\rho_{v,0}=1$ and $g_x,g_v<0$.
Normalization absorbs each nonzero diagonal coefficient into its gain and
does not change the interior operator.
A \texttt{Stencil} stores five weights in increasing offset order.
Construction requires finite weights and gains; it permits
nondecentralized stencils. Hypotheses of a particular approximation are
checked separately.

Nearest-neighbour coupling sets the two outer weights to zero.
With unit diagonal, specifying $\rho_{\bullet,1}$ determines
$\rho_{\bullet,-1}=-1-\rho_{\bullet,1}$.
The arguments \texttt{alpha\_x1} and \texttt{alpha\_v1} of
\texttt{nn\_stencil} denote these positive-offset weights.

\subsection{Circle, path and boundary conditions}

On a circle, $N$ counts all agents; indices are taken modulo $N$.
On a path, $N$ counts followers: agents are $0,\ldots,N$,
with agent $0$ the leader. Increasing index points towards the tail.
The standard leader experiment starts with $z_k(0)=0$,
$\dot z_k(0)=0$ for $k>0$, and $\dot z_0(0)=v_0$;
the leader follows $z_0(t)=v_0t$.
Its tail observable is $y_N=z_N-z_0$.
A circular velocity impulse is a different initial-value problem.

The endpoint rule is part of the operator. For the \texttt{drop} and
\texttt{regular} nearest-neighbour boundaries,
\begin{equation}\label{eq:tail}
 \ddot z_N=g_x b_x(z_N-z_{N-1})
          +g_v b_v(\dot z_N-\dot z_{N-1}).
\end{equation}
The \texttt{drop} boundary has $b_\bullet=-\rho_{\bullet,-1}$:
it discards missing neighbours and sets the diagonal to minus the
remaining off-diagonal sum. This is the variable-mass condition in
\cite[Definition 1]{transients}. The \texttt{regular} boundary has
$b_x=b_v=1$ and is defined only for nearest neighbours.
For next-nearest neighbours, \texttt{drop} applies its row rule at every
incomplete neighbourhood.

The next-nearest-neighbour experiments also implement the paper's fixed-mass
boundary, which clamps missing indices to the corresponding endpoint.
The general package boundary \texttt{fold} is a different rule:
it clamps missing indices except that, at the tail for $N\geq2$ and
$\rho_{-2}+\rho_{-1}\neq0$, it retains diagonal $\rho_0$ and multiplies
the predecessor weights by $1/|\rho_{-2}+\rho_{-1}|$.
That last rule need not preserve zero row sum for arbitrary input.
A record must identify which boundary was used.

\subsection{Finite operators and propagation}

For either geometry, let $L_x,L_v$ be the matrices implementing the stated
stencils and boundaries. In state order $(z,\dot z)$,
\begin{equation}\label{eq:matrix}
 \frac{\dd}{\dd t}\binom z{\dot z}
 =M\binom z{\dot z},\qquad
 M=\begin{pmatrix}0&I\\g_xL_x&g_vL_v\end{pmatrix}.
\end{equation}
A driven path has a zero acceleration row for its leader.
Removing its position and velocity coordinates gives the $2N$-dimensional
follower operator. The full path spectrum consists of that spectrum
together with the leader's double zero eigenvalue.

For a constant generator, the initial-value solution is
$u(t)=e^{tM}u(0)$. Pyflock evaluates a dense exponential, a sparse
exponential action, or a Fourier solution on the circle; an adaptive
DOP853 integrator provides a separate comparison.
These are floating-point evaluations of the finite system.
Neither an eigenvalue sensitivity heuristic nor agreement between two
numerical methods is a certified enclosure of the solution or spectrum.

\section{Fourier modes and local coefficients}

For the circle mode $z_k=e^{ik\theta}e^{\lambda t}$, put
\begin{equation}\label{eq:symbol}
 \nu_\bullet(\theta)=g_\bullet\sum_j\rho_{\bullet,j}e^{ij\theta}.
\end{equation}
Substitution in \eqref{eq:model} gives
\begin{equation}\label{eq:characteristic}
 \lambda^2-\nu_v(\theta)\lambda-\nu_x(\theta)=0.
\end{equation}
Thus the finite-circle spectrum, with algebraic multiplicity, is the union
of the two roots for $\theta=2\pi m/N$, $m=0,\ldots,N-1$.
This is an exact algebraic description; the returned numbers are numerical
root evaluations. Pointwise quadratic roots need not keep the same branch
label when followed around the circle.

Define moments
\[
 I_{\bullet,n}=g_\bullet\sum_j\rho_{\bullet,j}j^n,\qquad
 s^2=\frac{I_{v,1}^2}{4}+\frac{I_{x,2}}2.
\]
Assume decentralization, $I_{x,1}=0$, and $s>0$.
The two branches through the origin are locally analytic and have expansions
\begin{equation}\label{eq:expansion}
 \lambda_\pm(\theta)=-ic_\pm\theta-D_\pm\theta^2+O(\theta^3),
 \qquad c_\pm=-\frac{I_{v,1}}2\pm s.
\end{equation}
Positive velocity means increasing agent index, so phase velocity is
$-\operatorname{Im}\lambda/\theta$.

\begin{proposition}\label{prop:coefficients}
Under the preceding assumptions,
\begin{equation}\label{eq:diffusion}
 D_\pm=\frac{I_{v,2}}4
 \mp\frac{3I_{v,1}I_{v,2}+2I_{x,3}}{24s}.
\end{equation}
These are local coefficient identities; they do not require global
spectral stability.
\end{proposition}
\begin{proof}
Expand $\nu_\bullet=\sum_n I_{\bullet,n}(i\theta)^n/n!$.
The coefficients of $\theta^2$ and $i\theta^3$ in
\eqref{eq:characteristic} give
\[
 c^2+I_{v,1}c-\frac{I_{x,2}}2=0,\qquad
 D(2c+I_{v,1})=\frac{I_{v,2}c}{2}-\frac{I_{x,3}}6.
\]
Since $2c_\pm+I_{v,1}=\pm2s$, solving yields \eqref{eq:diffusion}.
\end{proof}

The function \texttt{taylor} obtains higher coefficients recursively from
the same characteristic equation, without numerical differentiation.
It also permits $s^2<0$, when first-order roots can have nonzero real
parts. Real-speed helpers require $s^2>0$.
The degenerate case $s=0$ and nonzero first position moment require
separate analysis; the local Taylor routine rejects them.
For $I_{x,1}\neq0$, fractional powers can occur.
A low-wavenumber expansion does not determine the remainder uniformly
over the circle or through arbitrarily many reflections.

Keeping the first two terms of \eqref{eq:expansion} formally yields
$u_t+cu_\xi=Du_{\xi\xi}$ for a wave envelope.
Gaussian spreading in this equation alone supplies no error estimate
for a reflected path wave or a tail extremum.

\section{Stability and transient measurements}

For a decentralized circle, coherent motion gives a two-dimensional invariant
subspace with a single Jordan block at zero. Asymptotic stability relative to
this motion requires the zero eigenvalue to have algebraic multiplicity
exactly two and all other eigenvalues in the open left half-plane.
Path stability refers to the follower operator with the leader removed.
Stability at a fixed size and bounds uniform as $N\to\infty$ are distinct
questions; a stable nonnormal matrix can have large transient responses.
A partial spectrum near zero cannot establish full spectral stability.

For symmetric position coupling, negative normalized gains, and
$c_-<0<c_+$, the path-wave comparison from
\cite[Theorem 3]{transients} uses
\begin{equation}\label{eq:prediction}
 \begin{split}
 P&=N\left(\frac1{c_+}-\frac1{c_-}\right),\\
 T_k&=\frac N{c_+}+(k-1)P,\qquad
 A_k=-\left(\frac{c_-}{c_+}\right)^{k-1}\frac{Nv_0}{c_+}.
 \end{split}
\end{equation}
Here $A_k$ are signed tail extrema; $2P$ is the full period.
The attenuation between the first and third extrema is
$A_3/A_1=(c_-/c_+)^2$, whereas consecutive absolute extrema have
ratio $|c_-/c_+|$.

The cited theorem assumes asymptotic stability, flock stability and
decay assumptions on Fourier coefficients of the second time derivative.
Its proof invokes the bulk-wave and boundary conjectures in that paper.
Equation \eqref{eq:prediction} is therefore used as a conditional
large-$N$ comparison for a fixed number of reflections.
The package's local stencil checks do not establish all those hypotheses.
The abrupt leader step and point impulses must also be distinguished from
smooth initial waves.

Measured extrema are extracted from sampled tail traces, with optional
local parabolic refinement at isolated extrema. Flat extrema are reported
at the midpoint of the plateau's time interval. Endpoints and monotone
plateaus are not counted as local extrema.
The convergence experiment estimates a full period from the first and third
extrema and reports relative errors using the predicted value as denominator.
It separately checks changes after coarsening the time samples.
Sampling sensitivity, variation between samples, and error of an asymptotic
formula measure different things.
A maximum on a finite recorded interval is not a maximum over all time.

\section{Heterogeneous arrays}

The module \texttt{heterogeneous} uses Lyons's positive gain magnitudes:
\begin{equation}\label{eq:heterogeneous}
 \ddot z=-\diag(g_{x,k})L_xz-\diag(g_{v,k})L_v\dot z,
 \qquad g_{x,k},g_{v,k}>0.
\end{equation}
Here $N$ counts all agents on either graph; the path has $N-1$ followers.
Position weights are $(-1/2,1,-1/2)$ and velocity weights are
$(-(1+\beta_v)/2,1,(\beta_v-1)/2)$.
The condition $|\beta_v|\leq1$ makes the neighbour velocity weights
nonpositive; it does not by itself establish heterogeneous stability.
Constant gains reproduce the homogeneous model with gains $-g_x,-g_v$.

For a periodically repeated cell, define
\[
 M_x=\left\langle\frac1{g_x}\right\rangle,\qquad
 V=\left\langle\frac{g_v}{g_x}\right\rangle.
\]
The first-order coefficients in \cite[Theorem 2.5.1]{lyons} are
\[
 c_{1,\pm}=-\frac{\beta_v V}{2M_x}
       \pm\sqrt{\frac{\beta_v^2V^2}{4M_x^2}+\frac1{2M_x}}.
\]
Lyons uses $\lambda=ic_1\theta+\cdots$; increasing-index velocities
are $-c_1$. His Bloch phase is per cell; for cell length $p$,
$\theta=\phi/p$. The package reports this convention explicitly.
The equivalent homogeneous gains $1/M_x$ and $V/M_x$ match these
first-order coefficients, without determining all finite spectra or
transient responses of differently ordered arrays.
Distinct-agent ensembles are finite numerical comparisons with their
gain arrays saved; they do not prove a random-medium limit.

\section{Computational evidence and source adaptations}

The named experiments retain their numerical arrays, parameters,
boundary and initial data, software versions, source snapshot and
checksums. Recomputing the arrays, redrawing saved arrays, and verifying
saved files are separate operations.
Checksums establish agreement with a saved record, not the truth of
a mathematical claim. Independent identities and numerical comparisons
in the tests supply further checks.

The coverage map records unavailable source samples and constructed
comparisons. In particular, the Type III stencil printed in
\cite[Figure 7]{nnn} has a nonzero row sum.
Changing its second coefficient from $-15/4$ to $+15/4$ restores
zero row sum and zero first position moment; the experiment records this
correction. The paper's fixed-mass velocity endpoint uses velocity
weights, correcting a printed position subscript.
New ensembles retain their selection rules and full coefficient arrays.

The matched-coupling studies~\cite{asymmetric,symmetry} refine
sensitive resonance neighbourhoods at higher precision; their reported
frequency maxima are refined samples over a stated band.
The force experiment~\cite{energy} checks the stated finite-system
energy and a constructed distributed-force comparison.
These computations do not certify supremum norms or asymptotic
scaling in array size.
The exploratory Gaussian amplitude correction and small-asymmetry
spectral helpers are identified as proposals in their API documentation;
they are not correctness criteria for the reproduced experiments.

\begin{thebibliography}{9}
\bibitem{signal}
C.~E. Cantos, J.~J.~P. Veerman and D.~K. Hammond,
\emph{Signal velocity in oscillator arrays},
Eur.\ Phys.\ J.\ Special Topics \textbf{225} (2016), 1115--1126.
\href{https://doi.org/10.1140/epjst/e2016-02665-0}{doi:10.1140/epjst/e2016-02665-0}.

\bibitem{transients}
C.~E. Cantos, D.~K. Hammond and J.~J.~P. Veerman,
\emph{Transients in the synchronization of asymmetrically coupled oscillator arrays},
Eur.\ Phys.\ J.\ Special Topics \textbf{225} (2016), 1199--1209.
\href{https://doi.org/10.1140/epjst/e2016-02658-y}{doi:10.1140/epjst/e2016-02658-y}.

\bibitem{nnn}
J.~Herbrych, A.~G. Chazirakis, N.~Christakis and J.~J.~P. Veerman,
\emph{Dynamics of locally coupled agents with next nearest neighbor interaction},
Differ.\ Equ.\ Dyn.\ Syst.\ \textbf{29} (2021), 487--509; online 2017.
\href{https://doi.org/10.1007/s12591-017-0377-3}{doi:10.1007/s12591-017-0377-3}.

\bibitem{lyons}
R.~G. Lyons, \emph{Linear Nearest Neighbor Flocks With All Distinct Agents},
PhD dissertation, Portland State University (2022).
\url{https://pdxscholar.library.pdx.edu/open_access_etds/5903/}.

\bibitem{asymmetric}
F.~M. Tangerman, J.~J.~P. Veerman and B.~D. Sto\v{s}i\'c,
\emph{Asymmetric decentralized flocks},
IEEE Trans.\ Automat.\ Control \textbf{57} (2012), 2844--2853.
\href{https://doi.org/10.1109/TAC.2012.2199150}{doi:10.1109/TAC.2012.2199150}.

\bibitem{symmetry}
J.~J.~P. Veerman,
\emph{Symmetry and stability of homogeneous flocks (a position paper)},
Portland State University working paper (2010).
\url{https://pdxscholar.library.pdx.edu/mth_fac/131/}.

\bibitem{energy}
I.~Herman, S.~Knorn and A.~Ahl\'en,
\emph{Disturbance scaling in bidirectional vehicle platoons with different
asymmetry in position and velocity coupling},
Automatica \textbf{82} (2017), 13--20.
\href{https://doi.org/10.1016/j.automatica.2017.04.010}{doi:10.1016/j.automatica.2017.04.010}.
\end{thebibliography}
\end{document}
